\documentclass[12pt]{article}
\usepackage[a4paper,margin=1in]{geometry}\usepackage[T1]{fontenc}
\usepackage{lmodern,microtype}\usepackage{amsmath,amssymb,amsthm,mathtools}
\usepackage{booktabs,array,enumitem,xcolor}\usepackage[numbers,sort&compress]{natbib}
\usepackage[colorlinks=true,linkcolor=blue!55!black,citecolor=blue!55!black,urlcolor=blue!55!black]{hyperref}
\linespread{1.07}
\newtheorem{theorem}{Theorem}[section]\newtheorem{proposition}[theorem]{Proposition}
\newtheorem{corollary}[theorem]{Corollary}
\theoremstyle{remark}\newtheorem{remark}[theorem]{Remark}
\newcommand{\ord}{\operatorname{ord}}

\title{Why a Fixed Quartic Cannot Supply a Growing Spectral Count\\
\large Bounded Degree, Repeated Support, and Local Finiteness}
\author{Liang Wang}\date{July 2026}
\begin{document}\maketitle
\begin{abstract}
RH-212--RH-218 extract substantial exact geometry from one physical edge
quartet.  This paper determines what even perfect control of that quartet
could contribute to an infinite spectral determinant.  The answer is sharply
limited.

First, a locally uniformly convergent family of monic polynomials of uniformly
bounded degree has a polynomial limit of bounded degree; it cannot acquire an
unbounded zero count.  Second, replacing a quartic $Q$ by powers $Q^N$ grows
algebraic degree only through multiplicity.  Its distinct support remains at
most four points, and as $N\to\infty$ the zero divisors cease to be locally
finite on every neighborhood of those points.  They therefore cannot
converge to the divisor of a nonzero holomorphic function on such a domain.

For the finest normalized physical quartet, powers through $N=64$ raise the
degree from four to 256 while the support remains exactly four.  The associated
height-counting functions have only the same finite jump locations with
growing jump sizes.  This behavior cannot produce a locally finite spectrum
with an extended counting law.

The conclusion is an exact route obstruction, not a negative result about
the quartet's local usefulness.  It remains a branch seed and finite factor,
but Gate A requires a rank-growing physical divisor that adds genuinely new
roots and controls omitted factors.  No $T\log T$ theorem or zeta statement is
claimed.
\end{abstract}

\section{The fixed-factor question}

The centered quartet has an exact shape manifold, a finite monotone axial
clock, and a classified degenerate boundary
\citep{WangRH213,WangRH214,WangRH216}.  Simple autonomous recurrence is not
identified, but suppose optimistically that a later argument controlled the
quartet at every scale.  Could this alone be the determinant sought in Gate
A?

An infinite spectral determinant must have a locally finite zero divisor:
every compact subset of its domain contains only finitely many zeros counted
with finite multiplicity, unless the function is identically zero.  A
four-root factor confronts this requirement in two possible ways:
\begin{enumerate}[leftmargin=2em]
\item keep degree four while taking a scale limit;
\item repeat or power the factor so algebraic degree grows.
\end{enumerate}
Both are insufficient.

\section{Bounded degree stays bounded}

\begin{theorem}[Locally uniform bounded-degree limit]\label{thm:degree}
Let $p_n$ be polynomials of degree at most $m$ that converge locally uniformly
on $\mathbb C$ to an entire function $f$.  Then $f$ is a polynomial of degree
at most $m$.  If every $p_n$ is monic of degree exactly $m$, then $f$ is monic
of degree $m$.
\end{theorem}
\begin{proof}
Local uniform convergence of holomorphic functions implies convergence of all
derivatives on compact subsets by Cauchy's integral formula
\citep{Conway1978}.  Since $p_n^{(m+1)}\equiv0$, one has
$f^{(m+1)}\equiv0$, so $f$ has degree at most $m$.  In the monic degree-$m$
case, $p_n^{(m)}\equiv m!$, hence $f^{(m)}\equiv m!$.
\end{proof}

\begin{corollary}\label{cor:four}
A locally uniform nonzero limit of monic quartics has exactly four zeros on
$\mathbb C$, counted with multiplicity.  It cannot yield an unbounded
spectral count.
\end{corollary}

On a proper domain, Rouch\'e's theorem gives the corresponding local
statement: if a compact contour avoids the zeros of the nonzero limit, the
number of enclosed zeros eventually stabilizes and is bounded by four
\citep{Conway1978}.

Thus even a rigorously proved limit
\begin{equation}
 Q_{\sigma}\longrightarrow(z^2-1)^2
\end{equation}
would produce one finite degenerate factor, not an infinite spectral
determinant.

\section{Powering grows multiplicity, not support}

Let
\begin{equation}
 Q(z)=\prod_{j=1}^{r}(z-\lambda_j)^{m_j},
 \qquad \sum_{j=1}^{r}m_j=4,qquad r\le4.
\end{equation}
Then
\begin{equation}\label{eq:powerdivisor}
 Q(z)^N=\prod_{j=1}^{r}(z-\lambda_j)^{Nm_j}.
\end{equation}

\begin{proposition}[Repeated-support obstruction]\label{prop:support}
For every $N\ge1$, the distinct zero support of $Q^N$ is the support of $Q$
and has cardinality at most four.  If $N\to\infty$, the divisors of $Q^N$ are
not locally finite on any domain containing a zero of $Q$.
\end{proposition}
\begin{proof}
Equation \eqref{eq:powerdivisor} proves the support statement.  If a compact
neighborhood contains $\lambda_j$, its divisor mass is at least $Nm_j$, which
diverges.  Local finiteness requires finite mass on compact sets.
\end{proof}

\begin{corollary}\label{cor:holomorphic}
No sequence $Q^{N_k}$ with $N_k\to\infty$ can have zero divisors converging
locally to the divisor of a nonzero holomorphic function on a domain
containing a zero of $Q$.
\end{corollary}

Multiplying by nonzero scalar normalizations does not change this divisor.
If the functions themselves converge locally uniformly, the only possible
way to accommodate unbounded compact zero order is an identically zero
limit, which is not a spectral determinant.

\section{Collapsing quartics do not repair local finiteness}

Suppose $Q_n$ tends to $(z^2-1)^2$ and one also chooses powers $N_n\to\infty$.
For any sufficiently small fixed neighborhoods of $+1$ and $-1$, all four
roots of $Q_n$ eventually lie in their union.  The divisor mass there is
$4N_n$ and diverges.  Root motion toward a compact degenerate pair therefore
does not create an extended locally finite support.

To obtain a genuine growing divisor, new roots must enter new locations in a
controlled manner.  Their multiplicities on every fixed compact set must
remain finite in the limiting object.

\section{Height counting profile}

For a finite root multiset $\Lambda$, define
\begin{equation}\label{eq:count}
 N_\Lambda(T)=\sum_{\lambda\in\Lambda}
 \mathbf 1_{\{|\operatorname{Im}\lambda|\le T\}},
\end{equation}
with multiplicity.  For $Q^N$,
\begin{equation}
 N_{Q^N}(T)=N\,N_Q(T).
\end{equation}
It has at most four distinct jump heights, independent of $N$.  Increasing
$N$ enlarges those jumps but does not distribute zeros over increasing
height.

This cannot approximate a locally finite extended counting law as a function
of $T$.  In particular it does not address the $T\log T$ scale required much
later in Gate C.  That macro-gate is not being proved or tested here; the
counting observation only explains why Gate A must not stop at one factor.

\section{Finite illustration}

At $\sigma=0.00125$ on the left channel, the centered quartet has
\begin{equation}
 u=0.718352,\qquad \eta=-0.093030.
\end{equation}
Its four roots are distinct.  We form powers with
\begin{equation}
 N=1,2,4,8,16,32,64.
\end{equation}

\begin{center}
\begin{tabular}{rrrr}
\toprule
$N$ & polynomial degree & distinct support & maximum root multiplicity\\
\midrule
$1$ & $4$ & $4$ & $1$\\
$8$ & $32$ & $4$ & $8$\\
$32$ & $128$ & $4$ & $32$\\
$64$ & $256$ & $4$ & $64$\\
\bottomrule
\end{tabular}
\end{center}

The archive also normalizes $Q(z)^N$ at a basepoint and samples five complex
locations.  As expected, modulus ratios are exponentially magnified or
suppressed with $N$; scalar normalization does not produce a stable divisor.
This numerical display illustrates Propositions~\ref{prop:support}; it is not
needed for their proof.

\section{What a rank-growing construction must add}

A viable next layer should contain divisors
\begin{equation}\label{eq:growing}
 D_{\sigma,k}(z)=\prod_{j=1}^{k}(z-\lambda_{\sigma,j}),
 \qquad k=k(\sigma)\longrightarrow\infty,
\end{equation}
with genuinely expanding support.  At minimum it must audit:
\begin{enumerate}[leftmargin=2em]
\item conjugate closure and branch selection as $k$ grows;
\item one global center and scale for the whole cloud;
\item tightness of normalized empirical root measures;
\item local uniform control of canonical products or logarithmic derivatives;
\item omitted-factor bounds connecting successive ranks;
\item dual-channel coherence.
\end{enumerate}

Adding roots is necessary, not sufficient.  Uncontrolled clouds can diverge
or converge to the zero function just as repeated factors do.

\section{Claim boundary}

The exact conclusion is a no-go for a fixed bounded-degree factor and its
pure repetitions as the final growing spectral divisor.  The quartet remains
useful as a local branch seed, a calibration factor, and a test of channel
universality.

No physical rank-growing divisor is constructed here.  Gate A remains open,
Gate C's counting law is untouched, and no arithmetic or zeta identification
is asserted.

\bibliographystyle{plainnat}\bibliography{references}\end{document}
